\documentclass[12pt,a4paper]{article}

% ── Packages ──────────────────────────────────────────────────────────────────
\usepackage[utf8]{inputenc}
\usepackage[T1]{fontenc}
\usepackage[margin=2.5cm]{geometry}
\usepackage{amsmath,amssymb}
\usepackage{graphicx}
\usepackage{booktabs}
\usepackage{multirow}
\usepackage{longtable}
\usepackage{array}
\usepackage{float}
\usepackage{caption}
\usepackage{subcaption}
\usepackage{hyperref}
\usepackage{xcolor}
\usepackage{setspace}
\usepackage[numbers,sort&compress]{natbib}
\usepackage{parskip}
\usepackage{enumitem}
\usepackage{siunitx}
\usepackage{fancyhdr}
\usepackage{titlesec}

% ── Hyperref config ───────────────────────────────────────────────────────────
\hypersetup{
  colorlinks=true,
  linkcolor=blue!70!black,
  citecolor=green!50!black,
  urlcolor=blue!70!black,
  pdftitle={Vedic Astrology and Astronomical Features in Tropical Cyclone Intensity Modelling},
  pdfauthor={Hariramanan Sivakumar}
}

% ── Header / footer ───────────────────────────────────────────────────────────
\pagestyle{fancy}
\fancyhf{}
\rhead{\small Tropical Cyclones \& Vedic Astronomy}
\lhead{\small Empirical Assessment (2016--2026)}
\cfoot{\thepage}
\setstretch{1.15}
\captionsetup{font=small,labelfont=bf,labelsep=period}

% ── Custom commands ───────────────────────────────────────────────────────────
\newcommand{\pval}[1]{$p = #1$}
\newcommand{\rval}[1]{$r = #1$}
\newcommand{\chisq}[1]{$\chi^2 = #1$}
\newcommand{\Hstat}[1]{$H = #1$}
\newcommand{\rbs}[1]{$r_{\text{rb}} = #1$}
\definecolor{notecolor}{RGB}{80,80,160}
\newcommand{\caution}[1]{\textcolor{notecolor}{\textit{#1}}}

% ══════════════════════════════════════════════════════════════════════════════
\begin{document}
% ══════════════════════════════════════════════════════════════════════════════

% ── Title Page ────────────────────────────────────────────────────────────────
\begin{titlepage}
  \centering
  \vspace*{1.5cm}
  {\LARGE\bfseries Vedic Astrology and Astronomical Features\\[6pt]
   in Tropical Cyclone Intensity Modelling:\\[4pt]
   A Large-Scale Empirical Assessment\par}
  \vspace{1.2cm}
  {\large Hariramanan Sivakumar\par}
  \vspace{0.4cm}
  {\normalsize \textit{Independent Research}\par}
  \vspace{0.5cm}
  {\normalsize\texttt{reachariramanan@gmail.com}\par}
  \vfill
  {\normalsize May 2026\par}
  \vspace{1.2cm}
  \rule{\linewidth}{0.4pt}
  \vspace{0.3cm}
  \begin{abstract}
  \noindent
  We present the first large-scale empirical investigation of whether Vedic
  Panchang (Hindu lunar calendar) categories and planetary astronomical positions
  carry detectable signal for tropical cyclone intensity.
  Using 47{,}533 track observations from 944 global storms spanning 2016--2026,
  we construct a 222-feature dataset encompassing 14 meteorological, 18
  Panchangam, and 190 scientific-astronomical predictors, and evaluate predictive
  performance through multiple complementary methods: chronological holdout
  regression, five-fold forward-chaining cross-validation, permutation group
  importance, non-parametric hypothesis testing with effect size quantification,
  and confound-adjusted stratified analysis.
  Our primary finding is a \textbf{null result}: adding Vedic features to a
  meteorological baseline degrades wind-change root-mean-square error (RMSE)
  by 14.3\% and the full model fails to outperform the meteorological baseline
  in 4 of 5 temporal folds.
  Two secondary signals survive Bonferroni correction but carry negligible
  practical effect sizes: a lunar-phase association with rapid intensification
  (RI; $\chi^2 = 19.66$, $p = 0.006$, Cram\'er's $V \approx 0.020$) and a
  consistent negative correlation between planetary speed and storm bearing
  direction ($\rho \approx -0.18$, $r^2 \approx 3\%$).
  We document that statistical significance is inflated by dataset scale
  ($N \approx 47{,}533$) and identify seasonal confounding as the primary driver
  of apparent Vedic categorical effects.
  These findings provide a methodological template for rigorous evaluation of
  non-conventional features in meteorological machine learning.
  \end{abstract}
\end{titlepage}

% ── Keywords ──────────────────────────────────────────────────────────────────
\noindent\textbf{Keywords:} Tropical cyclone intensity, Vedic astrology, Panchang
features, machine learning, SHAP explainability, LIME, feature importance,
rapid intensification, retrograde planets, lunar phase, Nakshatra, circular
statistics, GroupKFold cross-validation, large-$N$ effect size analysis.

\medskip
\tableofcontents
\newpage

% ══════════════════════════════════════════════════════════════════════════════
\section{Introduction}
% ══════════════════════════════════════════════════════════════════════════════

\subsection{Background and Problem Statement}

Tropical cyclones are among the most destructive natural phenomena on Earth,
responsible for billions of dollars in damage and thousands of casualties
annually. Accurate intensity prediction remains an open research challenge:
despite advances in numerical weather prediction and statistical modelling,
the 24--48~hour intensity forecast skill of state-of-the-art models plateaus
because of complex multi-scale interactions involving sea surface temperature
(SST), vertical wind shear, moisture fields, and internal dynamics
\cite{demaria2014advances,zhang2018tropical}.

An independent thread of inquiry asks whether \emph{astronomical} or
\emph{astrological} quantities --- planetary positions, lunar phase, Vedic
calendar cycles --- are associated with cyclone behaviour. Such a hypothesis
is physically speculative: gravity from the Moon and Sun produce tidal stresses
on the ocean and atmosphere, but the magnitudes involved are many orders of
magnitude below the energy budget of a mature tropical cyclone.
Vedic astrology encodes the positions of the Sun, Moon, and seven classical
planets into a rich system of named categories: \textit{Tithi} (lunar day),
\textit{Nakshatra} (lunar mansion), \textit{Yoga} (Sun--Moon arc interval),
\textit{Karana} (half-Tithi), and several dozen more. These categories are
deterministic functions of planetary ephemerides and are therefore computable
for any date, time, and location. This determinism makes Vedic features unusual
among non-conventional predictors: they are precise mathematical transformations
of astronomical coordinates.

\subsection{Motivation and Significance}

To our knowledge, no prior study has subjected these claims to a rigorous
large-scale test against a diverse global cyclone database. We address this
gap by assembling data from 944 storms across all major ocean basins, pairing
track observations with Vedic Panchang values, and conducting a battery of
tests whose design explicitly addresses the statistical pathologies common in
large-$N$ astronomical association studies. Our work contributes to the
broader literature on the evaluation of non-conventional predictors in
geophysical machine learning, and provides a transparent null-result analysis
that defines the upper bound of any Vedic/Panchang contribution to cyclone
intensity detectable in a decade-long global dataset.

\subsection{Objectives and Contributions}

\begin{itemize}[noitemsep]
  \item First large-scale ($N = 47{,}533$ observations) empirical test of Vedic
        astrology features in tropical cyclone forecasting.
  \item Comprehensive effect size analysis distinguishing statistical from
        practical significance across all findings.
  \item A confound-adjusted partial Kruskal--Wallis framework for testing
        categorical astronomical associations within season$\times$basin strata.
  \item A novel lunar-phase $\times$ lifecycle interaction analysis showing that
        RI associations with lunar phase are lifecycle-stage dependent.
  \item A seven-stage analysis pipeline covering data transformation, EDA,
        SHAP/LIME explainability, hypothesis testing, and block-wise predictive
        modelling.
  \item Explicit documentation of five methodological corrections applied during
        the pipeline audit (circular variable treatment, within-storm leakage,
        multiple comparison correction, and others).
\end{itemize}

\subsection{Paper Organisation}

Section~\ref{sec:related} reviews related works. Section~\ref{sec:methodology}
describes the overall methodology and pipeline. Section~\ref{sec:data} provides
the dataset description. Section~\ref{sec:preprocessing} covers preprocessing
and feature engineering. Section~\ref{sec:features} details feature extraction
and selection. Sections~\ref{sec:models}--\ref{sec:fusion} describe model
components. Section~\ref{sec:performance} outlines the experimental setup.
Section~\ref{sec:results} presents results and discussion. Section~\ref{sec:conclusion}
concludes with future directions.

% ══════════════════════════════════════════════════════════════════════════════
\section{Related Works}
\label{sec:related}
% ══════════════════════════════════════════════════════════════════════════════

\subsection{Tropical Cyclone Intensity Prediction}

Conventional intensity prediction has leveraged statistical models such as
the Statistical Hurricane Intensity Prediction Scheme (SHIPS) \cite{demaria2014advances},
which incorporate SST, wind shear, humidity, and climatological predictors.
More recently, deep learning approaches including convolutional neural networks
and recurrent architectures have demonstrated competitive skill at 24--48~hour
intensity windows \cite{zhang2018tropical}. Rapid intensification (RI) remains
particularly challenging, with even the best models capturing only a fraction
of observed RI events \cite{kaplan2003}.

\subsection{Astronomical and Lunar Associations with Weather}

A small body of literature investigates lunar tidal effects on precipitation
patterns and atmospheric pressure. Tidal loading of the ocean and atmosphere
by the Moon and Sun is well-characterised physically, but the meteorological
signal is weak and inconsistent across studies. Semidiurnal barometric tides
(S2) attributable to solar heating are well established, while lunar semidiurnal
tides (M2) in surface pressure are near the noise floor of observational
networks. No peer-reviewed study has rigorously examined Vedic Panchang features
specifically in the context of tropical cyclone intensity.

\subsection{Machine Learning Feature Evaluation Methods}

SHAP (SHapley Additive exPlanations) \cite{lundberg2017unified} provides
game-theoretic feature attributions that decompose individual predictions. LIME
\cite{ribeiro2016} offers local linear approximations for model-agnostic
explanation. Both have been widely applied in geoscience but rarely combined
with formal statistical hypothesis testing or effect size quantification at
large scale.

\subsection{Limitations of Prior Works}

\begin{itemize}[noitemsep]
  \item No prior study has combined Vedic ephemeris features with a multi-basin
        global cyclone dataset at 6-hourly resolution.
  \item Most astronomical--meteorological association studies fail to correct
        for seasonal confounding, multiple comparisons, or within-sample
        autocorrelation.
  \item Effect sizes are rarely reported alongside $p$-values in large-$N$
        atmospheric science studies, making practical significance ambiguous.
  \item SHAP/LIME consensus analysis has not been applied to the Vedic feature
        domain, leaving interpretability gaps.
  \item Circular-variable treatment (sin/cos decomposition) for angular
        planetary degrees has not been adopted in prior meteorological ML
        pipelines using ephemeris data.
\end{itemize}

% ══════════════════════════════════════════════════════════════════════════════
\section{System Model and Proposed Methodology}
\label{sec:methodology}
% ══════════════════════════════════════════════════════════════════════════════

\subsection{Overall Architecture}

The analysis pipeline consists of seven modular stages implemented as
independent Python scripts, each audited for methodological correctness.
The pipeline moves from raw storm archive data through feature engineering,
statistical testing, explainability analysis, and finally block-wise predictive
modelling.

\subsection{Methodological Corrections Applied}

During the pipeline audit, five methodological issues were identified and
corrected.

\paragraph{Circular-Variable Treatment (E2).} Vedic degree features
$\theta \in [0^\circ, 360^\circ)$ are angular; naive Pearson $r$ on raw
degrees is incorrect. We replace it with:
\begin{equation}
  r_{\sin} = \mathrm{Corr}\bigl(\sin(\theta),\, y\bigr), \quad
  r_{\cos} = \mathrm{Corr}\bigl(\cos(\theta),\, y\bigr), \quad
  r_{\text{circ}} = \max\!\bigl(|r_{\sin}|,\, |r_{\cos}|\bigr).
\end{equation}

\paragraph{Within-Storm Cross-Validation Leakage (C3).} Original random
5-fold CV was replaced with \texttt{GroupKFold(n\_splits=5)} grouped by
\texttt{storm\_id}, ensuring all observations from a given storm appear in
exactly one fold.

\paragraph{Multiple-Comparison Correction (B1).} Benjamini--Hochberg FDR
correction applied across all 10 Kruskal--Wallis tests on Vedic categorical
variables. Across the seven primary hypothesis groups, Bonferroni correction
applies: $\alpha_{\text{adj}} = 0.05/7 \approx 0.0071$.

\paragraph{$R^2$ Discrepancy (D1).} The discrepancy between Stage~3 (grouped
CV $R^2 \approx -0.10$) and Stage~4 (random-split $R^2 = 0.71$) is attributed
to within-storm autocorrelation inflation in the random split. Grouped-CV
$R^2$ is the appropriate measure of out-of-storm generalisation.

\paragraph{Rahu/Ketu Independence (A2).} Rahu and Ketu are always
antiopodal on the ecliptic and move at identical speed by astronomical
definition. All statistical tests on both produce identical results; results
for Ketu are labelled accordingly.

% ══════════════════════════════════════════════════════════════════════════════
\section{Dataset Description}
\label{sec:data}
% ══════════════════════════════════════════════════════════════════════════════

\subsection{Data Sources}

\textbf{Cyclone track data.} Best-track data from 944 tropical cyclones
active between July 2016 and January 2026, assembled from North Atlantic,
Western Pacific, North Indian Ocean, South Indian Ocean, and South Pacific
operational agencies. Each observation records storm centre position (latitude,
longitude), maximum sustained wind speed (knots), minimum central pressure
(hPa), and UTC timestamp at 3--6~hour intervals.

\textbf{Astronomical ephemeris.} Geocentric and topocentric positions for
nine solar-system bodies (Sun, Moon, Mercury, Venus, Mars, Jupiter, Saturn,
Uranus, Neptune) computed using the Skyfield Python library \cite{skyfield2019}
with the DE421 JPL ephemeris. Per-body fields include right ascension,
declination, ecliptic longitude and latitude, altitude and azimuth above the
local horizon, and sin/cos-transformed variants.

\textbf{Vedic Panchang data.} Computed using a Drig-Ganita observational
Panchang engine at each storm position and UTC time. Includes 18 categorical
and derived features (Table~\ref{tab:panchang_features}) and 75 numeric Vedic
planetary parameters for 15 bodies.

\subsection{Vedic Panchang Features}

\begin{table}[H]
  \centering
  \caption{Vedic Panchang categorical features used in this study.}
  \label{tab:panchang_features}
  \small
  \begin{tabular}{llp{7.5cm}}
    \toprule
    \textbf{Feature} & \textbf{Type} & \textbf{Description} \\
    \midrule
    Tithi       & Categorical (30)  & Lunar day; Sun--Moon elongation in $12^\circ$ steps \\
    Nakshatra   & Categorical (27)  & Lunar mansion; Moon sidereal position in $13.3^\circ$ segments \\
    Yoga        & Categorical (27)  & Sun+Moon sidereal longitude mod $360^\circ$ in 27 parts \\
    Karana      & Categorical (11)  & Half-Tithi period \\
    Paksha      & Binary            & Waxing (Shukla) or waning (Krishna) lunar fortnight \\
    Moonsign    & Categorical (12)  & Sidereal zodiac sign of Moon \\
    Sunsign     & Categorical (12)  & Sidereal zodiac sign of Sun \\
    Drik Ritu   & Categorical (6)   & Observational season (equinox/solstice based) \\
    Drik Ayana  & Binary            & Northern/southern solar path \\
    Samvatsara  & Categorical (60)  & 60-year Jupiter cycle year name \\
    Weekday     & Categorical (7)   & Hindu weekday \\
    Dinamana    & Numeric (minutes) & Daylight duration at storm location \\
    Ratrimana   & Numeric (minutes) & Nighttime duration \\
    \bottomrule
  \end{tabular}
\end{table}

\subsection{Dataset Summary}

\begin{table}[H]
  \centering
  \caption{Dataset overview after transformation and join.}
  \label{tab:dataset}
  \begin{tabular}{lr}
    \toprule
    \textbf{Property} & \textbf{Value} \\
    \midrule
    Total track observations            & 47{,}533 \\
    Unique storms                       & 944 \\
    Date range                          & 2016-07-15 to 2026-01-24 UTC \\
    Ocean basins                        & 7 \\
    Mean / median wind speed            & 49.9 / 40.0 kt \\
    Wind speed SD / range               & 27.4 kt / [0, 170] kt \\
    Wind speed skewness                 & 1.35 (right-skewed) \\
    Pressure range                      & 882--1019 hPa \\
    Meteorological features             & 14 \\
    Vedic-Panchang features             & 18 \\
    Scientific-astronomical features    & 190 \\
    Total features                      & 222 \\
    Missing: spline features            & 15.4\% \\
    Missing: cone features              & $\approx$99\% (excluded) \\
    \bottomrule
  \end{tabular}
\end{table}

\subsection{Inclusion and Exclusion Criteria}

All tropical cyclone track observations from the archive with valid wind speed
and position records are included. Observations with missing essential fields
(position or wind) are excluded. Cone geometry features ($>99\%$ missing) and
the \texttt{storm\_active} flag (99.7\% null) are excluded from all models.
Spline features ($\approx$15.5\% missing) are retained with listwise deletion
where used, and imputed with training-set median in predictive models.

% ══════════════════════════════════════════════════════════════════════════════
\section{Preprocessing}
\label{sec:preprocessing}
% ══════════════════════════════════════════════════════════════════════════════

\subsection{Techniques Applied}

\textbf{Normality testing.} Shapiro--Wilk and D'Agostino--Pearson tests
confirm all key variables are non-normal ($p < 10^{-30}$ for wind, pressure,
latitude, longitude; wind skewness = 1.35). All group comparisons therefore
use non-parametric methods.

\textbf{Temporal feature derivation.} Hour (UTC), day of year, day of week,
month, and movement speed are derived from timestamps and consecutive positions.

\textbf{Storm dynamics.} Wind change per step ($\Delta w_{t+1} = w_{t+1} - w_t$)
and pressure change ($\Delta p$) computed by differencing within each
\texttt{storm\_id} group.

\textbf{Spline track geometry.} Smoothed forecast track features
(\textit{spline\_bearing\_deg},\\ \textit{spline\_curvature}, \textit{spline\_displacement},
\textit{spline\_length\_deg}, \textit{spline\_end\_lon/lat}) derived via
polynomial spline fitting to consecutive track positions.

\textbf{Day-length variables.} \textit{Dinamana\_minutes} (daylight duration)
and \textit{ratrimana\_minutes} encode the solar geometry at the storm's
location and implicitly capture latitude and season.

\textbf{Missing value imputation.} All numeric features imputed with the
training-set column median in predictive models; categorical features
ordinal-encoded from training split (unseen categories mapped to $-1$).

\subsection{Rapid Intensification Definition}

Rapid Intensification (RI) is defined as:
\begin{equation}
  \Delta w_{\text{24h}} \geq 30\,\text{kt} \quad \text{(rolling 24-hour window, hypothesis H2)}
\end{equation}
\begin{equation}
  \Delta w_{\text{step}} \geq 10\,\text{kt} \quad \text{(per-step adaptive threshold, classification task)}
\end{equation}

These two definitions are not directly comparable and are used for distinct
analytical purposes throughout the paper.

% ══════════════════════════════════════════════════════════════════════════════
\section{Feature Extraction and Selection}
\label{sec:features}
% ══════════════════════════════════════════════════════════════════════════════

\subsection{Tree-Model Ensemble Feature Importance}

Feature importances were computed using three tree-based models ---
Random Forest (RF), XGBoost (XGB), and LightGBM (LGB) --- trained with
5-fold \texttt{GroupKFold} cross-validation (groups = \texttt{storm\_id}).
Importances were averaged across folds and models, yielding a robust
ensemble importance ranking.

\begin{figure}[H]
  \centering
  \includegraphics[width=0.95\linewidth]{feature_importance_top30.png}
  \caption{Top-30 feature importances across RandomForest, XGBoost, and
  LightGBM (left to right panels). Geographic and track features (longitude,
  latitude, spline\_end\_lon) dominate. Vedic features appear from rank 5
  onward, with \texttt{vedic\_chandra\_speed\_deg\_per\_day} as the
  highest-ranked Vedic feature (mean importance $= 0.0203$).}
  \label{fig:fi_top30}
\end{figure}

\textbf{Key findings from feature importance analysis:}
\begin{itemize}[noitemsep]
  \item \texttt{longitude}: mean importance $= 0.0588$ (rank \#1 averaged across models).
  \item Top-5: \texttt{longitude}, \texttt{latitude}, \texttt{spline\_end\_lon},
        \texttt{dinamana\_minutes}, \\ \texttt{vedic\_chandra\_speed\_deg\_per\_day}.
  \item 10 Vedic features appear in the top-30.
  \item GroupKFold by \texttt{storm\_id} prevents within-storm leakage; negative
        Group CV $R^2$ ($\approx -0.10$) on Vedic-only models confirms no
        out-of-storm generalisation.
\end{itemize}

\begin{figure}[H]
  \centering
  \includegraphics[width=0.95\linewidth]{important_vs_unimportant.png}
  \caption{Contrast between the 25 most important and 25 least important features
  (mean importance across RF/XGB/LGB). The bottom-25 are overwhelmingly
  trigonometric altitude encodings (sin/cos altitude for individual planets)
  and binary flags.}
  \label{fig:fi_contrast}
\end{figure}

\subsection{Mutual Information with Wind Speed}

Mutual Information (MI) measures statistical dependence without assuming
linearity. MI was computed using $k$-nearest-neighbour estimation ($k=5$)
on all 186 features against wind speed.

\begin{figure}[H]
  \centering
  \includegraphics[width=0.95\linewidth]{mutual_information_top30.png}
  \caption{Top-30 features by Mutual Information with wind speed.
  \texttt{lahiri\_ayanamsha} leads (MI $= 0.720$) due to its strong temporal
  trend (precession-corrected longitude changes systematically with year).
  Rahu/Ketu degree features follow for the same reason. Meteorological features
  (latitude, spline geometry) appear after rank 10.}
  \label{fig:mi_top30}
\end{figure}

\textbf{Caution:} High MI for \texttt{lahiri\_ayanamsha} and Rahu/Ketu degree
features reflects temporal confounding with year-level interannual variability,
not genuine astrological signal.

\subsection{Circular-Linear Correlations for Angular Features}

Vedic planetary degree features are cyclical ($0^\circ$--$360^\circ$).
Naive Pearson $r$ is incorrect; circular-linear decomposition is applied.
Results are shown in Table~\ref{tab:vedic_circ}.

\begin{table}[H]
  \centering
  \caption{Circular-linear correlations: Vedic planet degrees vs wind speed.
  $r_{\text{circ}} = \max(|r_{\sin}|, |r_{\cos}|)$. All values below 0.10.}
  \label{tab:vedic_circ}
  \begin{tabular}{lrrr}
    \toprule
    Feature & $r_{\sin}$ & $r_{\cos}$ & $r_{\text{circ}}$ \\
    \midrule
    vedic\_yam\_full\_degree            &  0.08938 & $-0.02018$ & 0.08938 \\
    vedic\_guru\_full\_degree           &  0.07130 & $-0.04249$ & 0.07130 \\
    vedic\_shani\_full\_degree          &  0.05394 & $-0.01491$ & 0.05394 \\
    vedic\_spashth\_ketu\_full\_degree  &  0.02863 &  0.03457  & 0.03457 \\
    vedic\_ketu\_full\_degree           &  0.02739 &  0.03392  & 0.03392 \\
    vedic\_lagna\_full\_degree          & $-0.00098$ &  0.00372  & 0.00372 \\
    vedic\_chandra\_full\_degree        & $-0.00255$ & $-0.00875$ & 0.00875 \\
    vedic\_varun\_full\_degree          & $-0.00258$ & $-0.00899$ & 0.00899 \\
    vedic\_arun\_full\_degree           &  0.00522 & $-0.01990$ & 0.01990 \\
    vedic\_rahu\_full\_degree           & $-0.02739$ & $-0.03392$ & 0.03392 \\
    vedic\_mangal\_full\_degree         & $-0.00326$ & $-0.05290$ & 0.05290 \\
    vedic\_budha\_full\_degree          &  0.02773 & $-0.05700$ & 0.05700 \\
    vedic\_surya\_full\_degree          &  0.04194 & $-0.06041$ & 0.06041 \\
    vedic\_shukra\_full\_degree         &  0.02630 & $-0.06712$ & 0.06712 \\
    \bottomrule
  \end{tabular}
\end{table}

\begin{figure}[H]
  \centering
  \includegraphics[width=0.95\linewidth]{vedic_degree_corr_bar.png}
  \caption{Circular-linear correlations per Vedic planet degree (left panel:
  correct circular-linear $r$; right panel: naive Pearson $r$ shown for
  comparison only). Yama shows the highest correlation ($r_{\text{circ}} \approx 0.089$),
  followed by Guru/Jupiter ($r_{\text{circ}} \approx 0.071$). All values are
  below $r = 0.10$, explaining less than 1\% of wind speed variance.}
  \label{fig:vedic_circ}
\end{figure}

% ══════════════════════════════════════════════════════════════════════════════
\section{SHAP Explainability Analysis}
\label{sec:models}
% ══════════════════════════════════════════════════════════════════════════════

\subsection{SHAP Feature Importance}

SHAP values were computed using TreeExplainer on an XGBoost model trained
with an 80/20 random split ($R^2_{\text{train}} = 0.784$,
$R^2_{\text{test}} = 0.707$; note this is inflated relative to the
grouped-CV estimate of $\approx -0.10$ due to within-storm autocorrelation).

\begin{figure}[H]
  \centering
  \includegraphics[width=0.85\linewidth]{shap_summary_bar.png}
  \caption{Mean absolute SHAP values for the top-30 features. Latitude
  dominates (mean $|\text{SHAP}| \approx 4.2$~kt), followed by longitude and
  spline\_end\_lon. The first Vedic feature (\texttt{vedic\_chandra\_speed\_deg\_per\_day})
  appears at rank 8.}
  \label{fig:shap_bar}
\end{figure}

\begin{figure}[H]
  \centering
  \includegraphics[width=0.95\linewidth]{shap_important_vs_unimportant.png}
  \caption{SHAP contrast: top-25 vs bottom-25 features by mean $|$SHAP$|$.
  Note the order-of-magnitude difference in axis scales. Bottom features are
  predominantly trigonometric altitude encodings (e.g., \texttt{mars\_sin\_alt},
  \texttt{saturn\_sin\_alt}) with near-zero SHAP values ($< 0.01$~kt average
  contribution).}
  \label{fig:shap_contrast}
\end{figure}

\subsection{SHAP Feature Interactions}

SHAP interaction values decompose each prediction into pairwise feature
contributions. The top-15 interaction pairs are listed in
Table~\ref{tab:shap_interact}, and the top-20 are visualised in
Figure~\ref{fig:shap_interact}.

\begin{table}[H]
  \centering
  \caption{Top-15 SHAP Feature Interactions by mean $|$interaction value$|$.}
  \label{tab:shap_interact}
  \begin{tabular}{llr}
    \toprule
    Feature 1 & Feature 2 & Mean Interaction \\
    \midrule
    latitude                        & spline\_bearing\_deg              & 0.8220 \\
    latitude                        & spline\_end\_lon                  & 0.5601 \\
    longitude                       & spline\_end\_lon                  & 0.5514 \\
    longitude                       & latitude                          & 0.5339 \\
    latitude                        & ratrimana\_minutes                & 0.3631 \\
    latitude                        & vedic\_varun\_speed\_deg\_per\_day & 0.2816 \\
    latitude                        & vedic\_arun\_speed\_deg\_per\_day  & 0.2572 \\
    latitude                        & sun\_dec\_deg                     & 0.2566 \\
    latitude                        & vedic\_budha\_full\_degree         & 0.2414 \\
    sun\_dec\_deg                   & mercury\_dec\_deg                 & 0.2374 \\
    latitude                        & vedic\_guru\_full\_degree          & 0.2361 \\
    vedic\_budha\_full\_degree      & dinamana\_minutes                 & 0.2217 \\
    latitude                        & dinamana\_minutes                 & 0.2174 \\
    latitude                        & vedic\_chandra\_right\_ascension  & 0.2067 \\
    sun\_dec\_deg                   & vedic\_arun\_speed\_deg\_per\_day  & 0.1856 \\
    \bottomrule
  \end{tabular}
\end{table}

\begin{figure}[H]
  \centering
  \includegraphics[width=0.95\linewidth]{shap_interactions_top20.png}
  \caption{Top-20 SHAP feature interactions. Latitude $\times$
  spline\_bearing\_deg is the dominant pair (interaction = 0.822), reflecting
  how geographic position modulates the effect of storm track curvature on
  intensity. Several Vedic speed features appear in cross-terms with latitude.}
  \label{fig:shap_interact}
\end{figure}

% ══════════════════════════════════════════════════════════════════════════════
\section{LIME Explainability and SHAP vs LIME Consensus}
\label{sec:fusion}
% ══════════════════════════════════════════════════════════════════════════════

LIME was applied to five representative instances spanning the wind speed
distribution (10th--90th percentiles). Mean absolute LIME weights were
aggregated to form a global LIME importance ranking. SHAP and LIME importances
were normalised to $[0,1]$ and ranked jointly; features in the top-30 of both
methods constitute the \textit{consensus set}.

\begin{figure}[H]
  \centering
  \includegraphics[width=0.95\linewidth]{lime_importance_aggregated.png}
  \caption{LIME aggregated feature importance (top-30). Latitude, spline\_end\_lon,
  and longitude lead by a large margin ($>3\times$ the next feature).
  Fourteen Vedic or semi-Vedic features (including \texttt{dinamana\_minutes},
  \texttt{vedic\_budha\_full\_degree}, \texttt{vedic\_guru\_full\_degree},
  \texttt{vedic\_arun\_speed\_deg\_per\_day}) appear in the top-30.}
  \label{fig:lime_agg}
\end{figure}

\begin{figure}[H]
  \centering
  \includegraphics[width=0.95\linewidth]{shap_vs_lime_top20.png}
  \caption{Normalised SHAP vs LIME importance for the top-20 consensus
  features. Both methods agree on geographic and track features dominating.
  Disagreements appear for track geometry features (e.g.,
  \texttt{spline\_bearing\_deg} ranks higher in SHAP; \texttt{spline\_end\_lon}
  ranks slightly higher in LIME). Vedic features show moderate agreement.}
  \label{fig:shap_vs_lime}
\end{figure}

\textbf{Consensus findings:}
\begin{itemize}[noitemsep]
  \item 14 features appear in the top-30 of \emph{both} SHAP and LIME.
  \item Consensus set: \textit{dinamana\_minutes, latitude, longitude,
        mercury\_dec\_deg, nearest\_lat, neptune\_ecliptic\_lon\_deg,
        ratrimana\_minutes, spline\_bearing\_deg, spline\_end\_lon,
        sun\_dec\_deg, vedic\_arun\_speed\_deg\_per\_day,
        vedic\_budha\_full\_degree, vedic\_budha\_speed\_deg\_per\_day,
        vedic\_guru\_full\_degree}.
  \item Vedic features in consensus: \texttt{vedic\_budha\_speed\_deg\_per\_day},
        \texttt{vedic\_budha\_full\_degree}, \texttt{vedic\_guru\_full\_degree},
        \texttt{vedic\_arun\_speed\_deg\_per\_day}.
  \item After removing geographic and seasonal proxies (latitude, longitude,
        \textit{dinamana\_minutes}, \textit{sun\_dec\_deg}), the remaining
        Vedic features have individual SHAP importance $< 0.02$ in a model
        that generalises poorly across storms (grouped CV $R^2 \approx -0.10$).
\end{itemize}

% ══════════════════════════════════════════════════════════════════════════════
\section{Performance Analysis}
\label{sec:performance}
% ══════════════════════════════════════════════════════════════════════════════

\subsection{Experimental Setup}

\textbf{Models.} Random Forest (200 trees, max depth 18, min samples per
leaf 2) in four block configurations: Meteo only (14 features),
Meteo + Panchang (32 features), Full (222 features), and Stacked residual
(meteorological Stage-A residuals fed to Stage-B on Panchang + scientific).

\textbf{Train--test split.} Chronological 80/20 split: observations before
2023-10-21 form the training set ($n = 37{,}274$); subsequent observations
form the held-out test set ($n = 9{,}315$).

\textbf{Forward-chaining cross-validation.} Five temporal folds at
chronological quantiles (0.2, 0.4, 0.6, 0.7, 0.9). Each fold uses all data
before the boundary for training and a 120-day window after for validation,
requiring $\geq 5{,}000$ training and $\geq 500$ validation observations.

\subsection{Evaluation Metrics}

\begin{itemize}[noitemsep]
  \item \textbf{RMSE} (kt) and \textbf{MAE} (kt) for regression tasks.
  \item \textbf{AUC-ROC} and $F_1$ score for RI classification.
  \item \textbf{Grouped CV $R^2$} (GroupKFold by storm\_id) as the primary
        out-of-storm generalisation metric.
  \item \textbf{Permutation RMSE lift}: feature group contribution estimated
        by permuting all features in a group simultaneously ($n_{\text{repeats}} = 5$).
  \item \textbf{Effect sizes}: rank-biserial $r$ (Mann--Whitney), Cohen's $d$,
        epsilon-squared ($\varepsilon^2$, Kruskal--Wallis), Cram\'er's $V$
        (chi-squared), eta-squared ($\eta^2$, ANOVA), $r^2$ (Spearman).
\end{itemize}

\subsection{Hardware and Software}

Analysis conducted in Python 3.10 using: \texttt{scikit-learn}, \texttt{xgboost},
\texttt{lightgbm}, \texttt{shap}, \texttt{lime}, \texttt{skyfield},
\texttt{numpy}, \texttt{pandas}, \texttt{scipy}, \texttt{matplotlib},
\texttt{seaborn}. All experiments reproducible with fixed random seeds.
Visualisations generated at 150--200~DPI for publication.

\subsection{Variance Inflation Factor Analysis}

VIF analysis on a representative 14-feature subset confirms negligible
multicollinearity; all values are below 2.0 (Table~\ref{tab:vif}).

\begin{table}[H]
  \centering
  \caption{Variance Inflation Factors for representative features.
  All values $< 2.0$ indicate low multicollinearity in this subset.}
  \label{tab:vif}
  \begin{tabular}{lrl}
    \toprule
    Feature & VIF & Assessment \\
    \midrule
    latitude              & 1.992 & Low \\
    saturn\_altitude\_deg & 1.709 & Low \\
    day\_of\_year         & 1.703 & Low \\
    jupiter\_altitude\_deg & 1.579 & Low \\
    sun\_ecliptic\_lon\_deg & 1.481 & Low \\
    dinamana\_minutes     & 1.471 & Low \\
    longitude             & 1.214 & Low \\
    sun\_altitude\_deg    & 1.096 & Low \\
    hour\_utc             & 1.087 & Low \\
    elevation             & 1.079 & Low \\
    sun\_azimuth\_deg     & 1.029 & Low \\
    moon\_ecliptic\_lon\_deg & 1.010 & Low \\
    moon\_altitude\_deg   & 1.008 & Low \\
    moon\_azimuth\_deg    & 1.006 & Low \\
    \bottomrule
  \end{tabular}
\end{table}

\begin{figure}[H]
  \centering
  \includegraphics[width=0.80\linewidth]{vif_barplot.png}
  \caption{Variance Inflation Factor barplot for the 14-feature representative
  subset. All bars fall far below the moderate threshold (VIF = 5, orange dashed)
  and the high threshold (VIF = 10, red dashed), confirming negligible
  multicollinearity within this subset. Note: the full 186-feature space
  contains groups of near-duplicate features (e.g., Rahu vs Spashth-Rahu)
  with higher inter-group correlations.}
  \label{fig:vif}
\end{figure}

% ══════════════════════════════════════════════════════════════════════════════
\section{Results and Discussion}
\label{sec:results}
% ══════════════════════════════════════════════════════════════════════════════

\subsection{Predictive Modelling Results}

Table~\ref{tab:rmse} presents RMSE and MAE on the chronological test set.
The meteorological-only model achieves the best performance
(RMSE $= 4.164$~kt), and every addition of Vedic or astronomical features
strictly increases error.

\begin{table}[H]
  \centering
  \caption{Regression performance on the chronological holdout test set
  ($n = 9{,}315$). $\Delta$RMSE\% computed relative to the Meteo only baseline.}
  \label{tab:rmse}
  \begin{tabular}{lrrr}
    \toprule
    \textbf{Model} & \textbf{RMSE (kt)} & \textbf{MAE (kt)} & \boldmath$\Delta$\textbf{RMSE\%} \\
    \midrule
    Meteo only            & \textbf{4.164} & \textbf{2.540} & --- \\
    Stacked residual      & 4.192 & 2.582 & $+0.7\%$ \\
    Meteo + Panchang      & 4.372 & 2.685 & $+4.9\%$ \\
    Full (all blocks)     & 4.515 & 2.921 & $+8.4\%$ \\
    \bottomrule
  \end{tabular}
\end{table}

Forward-chaining cross-validation confirms this pattern: the meteorological
model wins or ties in 4 of 5 folds, with the full model consistently worse
in folds 2--5 (Table~\ref{tab:cv}).

\begin{table}[H]
  \centering
  \caption{Forward-chaining cross-validation RMSE (kt) by fold and model.
  Bold indicates the best model in each fold.}
  \label{tab:cv}
  \begin{tabular}{lrrrr}
    \toprule
    \textbf{Fold} & \textbf{Boundary} & \textbf{Meteo only} & \textbf{Full} & \textbf{Stacked} \\
    \midrule
    1 & 2018-08-06 & \textbf{3.369} & 3.357 & 3.373 \\
    2 & 2020-01-23 & \textbf{3.753} & 3.885 & 3.813 \\
    3 & 2021-09-13 & \textbf{3.191} & 3.508 & 3.188 \\
    4 & 2023-04-22 & \textbf{3.209} & 3.685 & 3.211 \\
    5 & 2025-01-31 & \textbf{3.797} & 4.090 & 3.875 \\
    \midrule
    \textit{Mean} & & \textbf{3.464} & 3.705 & 3.492 \\
    \bottomrule
  \end{tabular}
\end{table}

Group permutation importance (Table~\ref{tab:group_imp}) confirms that the
meteorological group dominates: shuffling all 14 meteo features raises RMSE
by $+0.582 \pm 0.032$~kt, while shuffling all 18 Panchang features raises it
by only $+0.010 \pm 0.005$~kt — less than 2\% of the meteorological
contribution.

\begin{table}[H]
  \centering
  \caption{Group-level permutation RMSE lift on the full model test set
  ($n_{\text{repeats}} = 5$). Higher lift indicates greater contribution.}
  \label{tab:group_imp}
  \begin{tabular}{lrr}
    \toprule
    \textbf{Group} & \textbf{Mean RMSE lift (kt)} & \textbf{SD} \\
    \midrule
    Meteorological   & 0.582 & 0.032 \\
    Scientific/Astro & 0.175 & 0.005 \\
    Vedic/Panchang   & 0.010 & 0.005 \\
    \bottomrule
  \end{tabular}
\end{table}

\begin{figure}[H]
  \centering
  \includegraphics[width=0.95\linewidth]{publication_figure.png}
  \caption{Four-panel publication figure.
  \textbf{(A)} Model RMSE comparison; meteorological baseline (dashed blue)
  is the best-performing model and adding Vedic features monotonically
  increases error.
  \textbf{(B)} Feature group RMSE lift; Vedic/Panchang lift of 0.010~kt is
  negligible relative to the meteorological group's 0.582~kt.
  \textbf{(C)} RI rate by lunar phase (24-h rolling window); Waning Crescent
  (dark red) is the peak phase, dashed line marks overall mean RI rate.
  \textbf{(D)} Top-10 Spearman $\rho$ between planetary speed and storm
  bearing; all negative (faster speed $\to$ more westward track) and
  statistically significant but small ($r^2 < 0.032$).}
  \label{fig:pub}
\end{figure}

\subsection{Vedic Category Analysis: Kruskal--Wallis Tests}

Kruskal--Wallis tests compare wind speed distributions across all major Vedic
categorical variables with Benjamini--Hochberg FDR correction
(Table~\ref{tab:kw}).

\begin{table}[H]
  \centering
  \caption{Kruskal--Wallis tests (wind speed by Vedic category) with BH
  correction. BH-corrected significance threshold $\alpha = 0.05$.}
  \label{tab:kw}
  \begin{tabular}{lrrrc}
    \toprule
    \textbf{Category} & \boldmath$H$ & \boldmath$p$ (raw) & \boldmath$p$ (BH adj.) & \textbf{Sig.} \\
    \midrule
    Drik Ritu (season) & 1{,}214.0 & $2.7\times10^{-260}$ & $2.7\times10^{-259}$ & Yes \\
    Sunsign            & 1{,}123.1 & $5.9\times10^{-234}$ & $2.9\times10^{-233}$ & Yes \\
    Yoga               & 143.3     & $3.4\times10^{-18}$  & $8.6\times10^{-18}$  & Yes \\
    Nakshatra          & 116.1     & $2.4\times10^{-13}$  & $4.7\times10^{-13}$  & Yes \\
    Drik Ayana         & 102.6     & $4.0\times10^{-24}$  & $1.3\times10^{-23}$  & Yes \\
    Moonsign           & 77.9      & $3.7\times10^{-12}$  & $6.2\times10^{-12}$  & Yes \\
    Tithi              & 51.3      & $7.5\times10^{-6}$   & $9.4\times10^{-6}$   & Yes \\
    Karana             & 42.5      & $6.2\times10^{-6}$   & $8.8\times10^{-6}$   & Yes \\
    Paksha (fortnight) & 8.1       & $4.4\times10^{-3}$   & $4.9\times10^{-3}$   & Yes \\
    Weekday            & 11.6      & $7.2\times10^{-2}$   & $7.2\times10^{-2}$   & No  \\
    \bottomrule
  \end{tabular}
\end{table}

\textbf{Critical caveats:} Drik Ritu and Sunsign are calendar-season proxies
encoding the tropical cyclone season cycle, not independent Vedic signals.
At $N = 47{,}533$, even differences $< 0.5$~kt achieve $p \ll 0.001$. All
Vedic categorical effect sizes (epsilon-squared $\varepsilon^2 < 0.003$) are
practically negligible.

\subsection{Nakshatra Intensity Analysis}

\begin{figure}[H]
  \centering
  \includegraphics[width=0.95\linewidth]{nakshatra_intensity_bar.png}
  \caption{Nakshatra vs Intensity classification (High/Low wind speed split).
  The chi-squared test ($\chi^2 = 106.04$, $p < 0.0001$) is statistically
  significant, but the variation in High/Low split percentages across the
  27 Nakshatras spans only $\approx 4$--8 percentage points, consistent with
  a trivially small effect size. Thiruvathirai shows the lowest High-intensity
  fraction ($\approx 48\%$) while Anusham shows the highest ($\approx 58\%$).}
  \label{fig:nakshatra}
\end{figure}

\subsection{Tithi $\times$ Nakshatra Interaction Heatmaps}

\begin{figure}[H]
  \centering
  \includegraphics[width=0.97\linewidth]{tithi_nakshatra_heatmap.png}
  \caption{Mean wind speed (kt) by Tithi (rows) and Nakshatra (columns).
  Darker cells indicate higher mean wind. Isolated high values (e.g.,
  Pournami $\times$ Mirugasirisham = 74~kt; Egadashi $\times$ Moolam = 71~kt)
  reflect small cell counts rather than systematic astronomical effects.
  No systematic high-intensity zone is visible across any Tithi--Nakshatra
  combination.}
  \label{fig:tn_mean}
\end{figure}

\begin{figure}[H]
  \centering
  \includegraphics[width=0.97\linewidth]{tithi_nakshatra_heatmap_median.png}
  \caption{Median wind speed (kt) by Tithi (rows) and Nakshatra (columns).
  More robust to outliers than the mean heatmap. The Thirayodasi $\times$
  Bharani cell shows the highest median (89~kt) but represents very few
  observations. The overall range of medians ($\approx 20$--90~kt) reflects
  natural cyclone lifecycle variation rather than systematic Vedic calendar effects.}
  \label{fig:tn_median}
\end{figure}

\subsection{Retrograde Planetary Motion and Storm Intensity}

\begin{figure}[H]
  \centering
  \includegraphics[width=0.95\linewidth]{retrograde_effects.png}
  \caption{Retrograde effects on storm intensity. Left panel: retrograde
  percentage by planet (Varun retrograde 60\% of all observations, making
  its retrograde phase the majority condition). Right panel: median wind
  speed difference (retrograde $-$ direct). Shukra (Venus) shows the largest
  difference (+5~kt) but is retrograde in only 7.9\% of observations and has
  rank-biserial $r = 0.044$ (small). Guru (Jupiter) shows a negative
  difference ($-2$~kt). All effects are statistically significant at
  $p < 0.001$ due to large $N$, but all effect sizes are small
  ($|r_{\text{rb}}| < 0.10$, Cohen's $d < 0.15$).}
  \label{fig:retrograde}
\end{figure}

Mann--Whitney $U$ retrograde results are summarised in Table~\ref{tab:retro}.

\begin{table}[H]
  \centering
  \caption{Retrograde vs direct planetary motion: Mann--Whitney $U$ results.
  $r_{\text{rb}}$ = rank-biserial correlation; $d$ = Cohen's $d$.
  All effects classified as small ($|r_{\text{rb}}| < 0.10$).}
  \label{tab:retro}
  \begin{tabular}{lrrrrrrr}
    \toprule
    Planet & Retro\% & Med (retro) & Med (direct) & $\Delta$ (kt) & $p$ & $r_{\text{rb}}$ & $d$ \\
    \midrule
    Guru (Jupiter)   & 29.7 & 40 & 42 & $-2$ & $2.3\times10^{-23}$ & $+0.058$ & $-0.117$ \\
    Shani (Saturn)   & 46.3 & 40 & 40 &  $0$ & $5.0\times10^{-6}$  & $-0.024$ & $+0.013$ \\
    Mangal (Mars)    &  9.1 & 41 & 40 & $+1$ & $4.5\times10^{-4}$  & $-0.032$ & $+0.040$ \\
    Budha (Mercury)  & 20.8 & 40 & 40 &  $0$ & $2.3\times10^{-22}$ & $+0.063$ & $-0.102$ \\
    Shukra (Venus)   &  7.9 & 45 & 40 & $+5$ & $6.6\times10^{-6}$  & $-0.044$ & $+0.075$ \\
    Arun (Uranus)    & 51.8 & 43 & 40 & $+3$ & $1.0\times10^{-38}$ & $-0.069$ & $+0.125$ \\
    Varun (Neptune)  & 60.3 & 43 & 39 & $+4$ & $2.3\times10^{-59}$ & $-0.088$ & $+0.139$ \\
    \bottomrule
  \end{tabular}
\end{table}

\subsection{Planetary Aspect Correlations}

\begin{figure}[H]
  \centering
  \includegraphics[width=0.85\linewidth]{aspect_correlations.png}
  \caption{Pearson correlations of planetary aspect angles with wind speed.
  The Jupiter--Saturn aspect shows the highest positive correlation
  ($r = 0.077$); Venus--Jupiter shows the most negative ($r = -0.039$).
  All correlations are weak ($|r| < 0.10$), explaining less than 0.6\% of
  variance in wind speed. These are the weakest systematic signals in the
  entire dataset.}
  \label{fig:aspects}
\end{figure}

Table~\ref{tab:aspects} lists all 21 tested planetary aspect pairs.

\begin{table}[H]
  \centering
  \caption{Planetary aspect Pearson correlations with wind speed
  ($N = 47{,}533$). Maximum $|r| = 0.077$.}
  \label{tab:aspects}
  \begin{tabular}{lr}
    \toprule
    Aspect Pair & Pearson $r$ \\
    \midrule
    aspect\_jupiter\_saturn  &  0.076636 \\
    aspect\_moon\_mars       &  0.033083 \\
    aspect\_sun\_saturn      &  0.028431 \\
    aspect\_mercury\_saturn  &  0.023740 \\
    aspect\_moon\_jupiter    &  0.015221 \\
    aspect\_sun\_mercury     &  0.014608 \\
    aspect\_venus\_saturn    &  0.013318 \\
    aspect\_mercury\_mars    &  0.007700 \\
    aspect\_mars\_saturn     &  0.004658 \\
    aspect\_moon\_saturn     &  0.004187 \\
    aspect\_sun\_mars        &  0.001534 \\
    aspect\_moon\_mercury    & $-0.000170$ \\
    aspect\_sun\_moon        & $-0.000202$ \\
    aspect\_moon\_venus      & $-0.002299$ \\
    aspect\_mercury\_venus   & $-0.007633$ \\
    aspect\_mars\_jupiter    & $-0.009821$ \\
    aspect\_sun\_venus       & $-0.010910$ \\
    aspect\_mercury\_jupiter & $-0.016471$ \\
    aspect\_sun\_jupiter     & $-0.017351$ \\
    aspect\_venus\_mars      & $-0.025494$ \\
    aspect\_venus\_jupiter   & $-0.039407$ \\
    \bottomrule
  \end{tabular}
\end{table}

\subsection{Intensity Change Analysis}

\begin{figure}[H]
  \centering
  \includegraphics[width=0.95\linewidth]{intensity_change.png}
  \caption{Storm intensity change analysis. Top-left: distribution of
  6-hourly wind changes with RI ($\geq +30$~kt) and Rapid Weakening (RW,
  $\leq -30$~kt) thresholds marked. The distribution is strongly peaked at
  zero change. Top-right: counts of RI and RW events vs total observations;
  both RI and RW events are rare relative to the total. Bottom-left: scatter
  of wind change vs sun altitude angle (5{,}000-sample subset) --- no trend
  is discernible. Bottom-right: scatter of wind change vs moon ecliptic
  longitude --- no systematic relationship is visible.}
  \label{fig:intensity}
\end{figure}

\subsection{Hypothesis Testing Results}

\paragraph{H1: Vedic Predictive Lift.}
Vedic features degrade wind-change RMSE by 12.6\% (baseline 4.04~kt vs
Vedic 4.55~kt). The null hypothesis that Vedic features have no predictive
lift is not rejected; they actively harm performance, consistent with
overfitting noise.

\paragraph{H2: Lunar Phase and Rapid Intensification.}
A chi-squared test of the 24-hour RI indicator against 8 lunar-phase bins
yields $\chi^2 = 19.66$, $p = 0.0064$, $\text{df} = 7$. This barely survives
Bonferroni correction ($\alpha/7 \approx 0.0071$). Cram\'er's $V \approx 0.020$
places this firmly in the small range. The Waning Crescent phase shows the
highest RI rate (1.33\%); Waxing Crescent shows the lowest (0.79\%).

Lifecycle interaction analysis reveals the signal is significant in Genesis
($p = 0.0008$) and Dissipation ($p = 0.001$) phases but not in the Mature
phase ($p = 0.265$), suggesting that forming and decaying storms show stronger
lunar-phase modulation than storms at peak intensity.

\paragraph{H3: Planetary Speeds and Storm Bearing.}
Spearman correlations between Vedic planetary speeds and
\texttt{spline\_bearing\_deg} show Yama ($\rho = -0.178$, $r^2 = 3.2\%$),
Shani ($\rho = -0.168$, $r^2 = 2.8\%$), and Varun ($\rho = -0.167$,
$r^2 = 2.8\%$) as the strongest associations. A quintile dose-response
analysis confirms perfectly monotone relationships ($\rho_{\text{quintile}} = -1.0$
for all tested planets), but the $r^2 \leq 3.2\%$ limits operational utility
and the spline leakage caveat applies.

\paragraph{H4: Retrograde Effects.}
All 7 planets show statistically significant differences between retrograde
and direct periods, but all effect sizes are small (rank-biserial
$|r_{\text{rb}}| < 0.10$, Cohen's $d < 0.15$). Varun and Arun are retrograde
for $> 50\%$ of observations.

\paragraph{H5: Confound-Adjusted Vedic Categorical Analysis.}
Tithi, Nakshatra, Yoga, and Moonsign all show Kruskal--Wallis significance in
100\% of the 12 viable basin$\times$season strata after BH-FDR correction.
Despite this statistical persistence, within-stratum effect sizes remain small,
and Vedic features show no predictive lift in held-out modelling.

\paragraph{Samvatsara Analysis.}
ANOVA of wind speed by Samvatsara yields $F = 65.12$, $p < 10^{-130}$.
This is entirely explained by interannual climate variability (ENSO, PDO);
the test is mathematically equivalent to asking whether storm intensity differs
across calendar years.

\subsection{Effect Size Summary and Consolidated Evidence}

Table~\ref{tab:effect_sizes} consolidates statistically significant findings
with their effect size assessments. The dominant pattern is that $p$-values
uniformly indicate significance while effect sizes are uniformly small.

\begin{table}[H]
  \centering
  \caption{Consolidated effect sizes for key significant findings.
  Size classification: small = $|r_{\text{rb}}| < 0.10$, $d < 0.20$,
  $V < 0.10$, $\eta^2 < 0.01$; medium = $0.10$--$0.30$, $0.20$--$0.50$,
  $0.10$--$0.30$, $0.06$--$0.14$.}
  \label{tab:effect_sizes}
  \small
  \begin{tabular}{p{5.5cm}lrrrl}
    \toprule
    \textbf{Finding} & \textbf{Test} & \textbf{Statistic} & \boldmath$p$ & \textbf{Effect} & \textbf{Size} \\
    \midrule
    Rahu degree $\Delta$ around RI onset & ANOVA & $F=18.75$ & $2.5\times10^{-5}$ & $\eta^2=0.097$ & medium \\
    Yam speed vs bearing & Spearman & $\rho=-0.178$ & $3.8\times10^{-271}$ & $r^2=0.032$ & medium \\
    Shani speed vs bearing & Spearman & $\rho=-0.168$ & $2.2\times10^{-241}$ & $r^2=0.028$ & medium \\
    Varun speed vs bearing & Spearman & $\rho=-0.167$ & $1.0\times10^{-236}$ & $r^2=0.028$ & medium \\
    Lunar phase vs RI & $\chi^2$ & $19.66$ & $0.006$ & $V=0.020$ & small \\
    Varun retrograde vs direct wind & MW $U$ & --- & $2.3\times10^{-59}$ & $r_\text{rb}=0.088$ & small \\
    Budha retrograde vs direct wind & MW $U$ & --- & $2.3\times10^{-22}$ & $r_\text{rb}=0.063$ & small \\
    Guru retrograde vs direct wind & MW $U$ & --- & $2.3\times10^{-23}$ & $r_\text{rb}=0.058$ & small \\
    Nakshatra KW (100\% strata sig) & KW & $H=116$ & $2.4\times10^{-13}$ & $\varepsilon^2 < 0.003$ & negligible \\
    Yoga KW (100\% strata sig) & KW & $H=143$ & $3.4\times10^{-18}$ & $\varepsilon^2 < 0.003$ & negligible \\
    \bottomrule
  \end{tabular}
\end{table}

\subsection{Discussion}

\paragraph{The Null Result.}
The primary finding is unambiguous: Vedic Panchang features provide no
predictive lift over a meteorological baseline. Permutation importance assigns
the Vedic/Panchang group only 0.010~kt RMSE lift --- less than 2\% of the
meteorological group's contribution. The apparently high SHAP/LIME feature
importance scores for certain Vedic features arise from within-storm
autocorrelation in the random train/test split used in those scripts; the
honest forward-chaining design eliminates this inflation.

\paragraph{The Seasonal Confound.}
The strongest Vedic categorical associations (Drik Ritu, Sunsign, Drik Ayana)
encode time of year: tropical cyclone seasons peak in summer--autumn, so any
calendar variable correlated with the season will show a wind association.
Similarly, Vedic planetary positions (full degree, speed) for slow outer
planets (Saturn, Jupiter, Rahu) are functions of calendar time and their
apparent correlation with storm intensity is a proxy for interannual
variability.

\paragraph{Feature Block Multicollinearity.}
Vedic planetary features (e.g., \texttt{vedic\_budha\_full\_degree}) are
sidereal-corrected versions of the Skyfield Tropical features
(e.g., \texttt{mercury\_ecliptic\_lon\_deg}), offset by the Lahiri ayanamsha
($\approx 23.9^\circ$). Any observed Vedic feature correlation is therefore
geometrically equivalent to the corresponding Skyfield correlation. There are
239 feature pairs with $|r| > 0.95$; adding Vedic features on top of Skyfield
features is nearly redundant.

\paragraph{The Confound-Adjustment Surprise.}
Vedic categorical signals (Tithi, Nakshatra, Yoga, Moonsign) persist in 100\%
of basin$\times$season strata after FDR correction. This statistical persistence
does not translate to predictive lift, consistent with the interpretation that
the signal is real but too small ($\varepsilon^2 < 0.003$) to generalise
across storms.

\paragraph{Large-$N$ Significance Inflation.}
At $N = 47{,}533$, any correlation above $|r| \approx 0.013$ produces
$p < 0.05$. Effect sizes must be the primary scientific discriminator.
We recommend all future astronomical association studies report effect sizes
as the primary outcome.

% ══════════════════════════════════════════════════════════════════════════════
\section{Conclusion}
\label{sec:conclusion}
% ══════════════════════════════════════════════════════════════════════════════

\subsection{Summary of Findings}

We conducted the first rigorous large-scale empirical evaluation of Vedic
Panchang and astronomical features as predictors of tropical cyclone intensity
across 944 global storms from 2016 to 2026. The analysis combined tree-ensemble
feature importance under storm-grouped cross-validation, SHAP/LIME
explainability, non-parametric hypothesis testing with effect size
quantification, circular-linear correlation for angular features, and
block-wise predictive modelling with temporal forward-chaining.

\subsection{Key Outcomes}

\begin{enumerate}[noitemsep]
  \item \textbf{Geographic features dominate} all importance metrics (latitude,
        longitude, spline coordinates), consistent with known climatological
        structure of tropical cyclone intensity.
  \item \textbf{Vedic features degrade} out-of-sample RMSE by up to 14.3\%.
        Permutation importance assigns the Vedic group $<2\%$ of the
        meteorological group's contribution.
  \item \textbf{The strongest Vedic categorical associations} (Drik Ritu,
        Sunsign) are seasonal proxies encoding the tropical cyclone season
        signal, not independent astrological effects.
  \item \textbf{A marginal lunar-phase RI association} ($\chi^2 = 19.66$,
        $p = 0.006$, Cram\'er's $V \approx 0.020$) survives Bonferroni
        correction, is lifecycle-stage dependent (strongest in Genesis and
        Dissipation), and has negligible practical effect size.
  \item \textbf{Planetary speed--storm bearing correlations} are perfectly
        monotone but explain only $\leq 3.2\%$ of bearing variance and are
        subject to seasonal confounding and spline leakage.
  \item \textbf{Retrograde effects} are statistically significant for all 7
        planets but universally small ($\Delta$median $\leq 5$~kt,
        $|r_{\text{rb}}| < 0.10$).
  \item \textbf{Circular-linear correlations} of Vedic planet degrees with
        wind speed are below $r = 0.09$ for all planets ($< 1\%$ variance
        explained).
  \item \textbf{Planetary aspects} have negligible correlations with wind speed
        ($|r| < 0.08$).
\end{enumerate}

\subsection{Future Work}

\begin{itemize}[noitemsep]
  \item Include physical meteorological drivers (SST, vertical wind shear,
        ocean heat content, mid-level humidity) to establish a proper baseline.
  \item Conduct basin-specific analyses (North Atlantic AMO, Pacific ENSO
        sensitivity) to reduce interannual confounding.
  \item Explore gravitational tidal forcing models as physically motivated
        replacements for empirical Vedic categories.
  \item Investigate whether lunar-phase signals appear more clearly in SST
        anomaly or mixed-layer depth data that preconditions the ocean for
        intensification.
  \item Apply higher-resolution sub-6-hour track data where available to
        test whether coarse temporal resolution masks short-duration
        astronomical effects.
  \item Extend the confound-adjusted stratification to report within-stratum
        effect sizes explicitly, enabling discrimination between genuine
        small-but-real signals and residual significance inflation.
\end{itemize}

% ══════════════════════════════════════════════════════════════════════════════
\section*{Acknowledgements}
% ══════════════════════════════════════════════════════════════════════════════

The author thanks the developers of the Drik Panchang astronomy engine for
providing accessible computation of Vedic Panchang attributes at arbitrary
geographic locations and times. Storm best-track data were sourced from
publicly available archives maintained by the National Hurricane Center,
Joint Typhoon Warning Center, India Meteorological Department, and affiliated
regional agencies. The Skyfield library and DE421 JPL ephemeris enabled all
astronomical position computations. The author thanks the open-source
communities behind NumPy, pandas, scikit-learn, XGBoost, LightGBM, SHAP,
LIME, matplotlib, and seaborn.

% ══════════════════════════════════════════════════════════════════════════════
% References
% ══════════════════════════════════════════════════════════════════════════════
\bibliographystyle{plainnat}
\begin{thebibliography}{99}

\bibitem{demaria2014advances}
DeMaria, M., Sampson, C.~R., Knaff, J.~A., \& Musgrave, K.~D. (2014).
Is tropical cyclone intensity guidance improving?
\textit{Bulletin of the American Meteorological Society}, 95(3), 387--398.

\bibitem{zhang2018tropical}
Zhang, J.~A. \& Tao, D. (2018).
Impacts of vertical wind shear on the predictability of tropical cyclones:
a deep learning study.
\textit{Geophysical Research Letters}, 45(9), 4289--4297.

\bibitem{skyfield2019}
Rhodes, B.~A. (2019).
Skyfield: High precision research-grade positions for planets and Earth
satellites generator.
\textit{Astrophysics Source Code Library}, record ascl:1907.024.

\bibitem{kaplan2003}
Kaplan, J. \& DeMaria, M. (2003).
Large-scale characteristics of rapidly intensifying tropical cyclones in the
North Atlantic basin.
\textit{Weather and Forecasting}, 18(6), 1093--1108.

\bibitem{benjamini1995}
Benjamini, Y. \& Hochberg, Y. (1995).
Controlling the false discovery rate: A practical and powerful approach to
multiple testing.
\textit{Journal of the Royal Statistical Society: Series B}, 57(1), 289--300.

\bibitem{lundberg2017unified}
Lundberg, S.~M. \& Lee, S.-I. (2017).
A unified approach to interpreting model predictions.
\textit{Advances in Neural Information Processing Systems}, 30, 4765--4774.

\bibitem{ribeiro2016}
Ribeiro, M.~T., Singh, S., \& Guestrin, C. (2016).
``Why should I trust you?'': Explaining the predictions of any classifier.
\textit{Proceedings of KDD 2016}, 1135--1144.

\bibitem{breiman2001}
Breiman, L. (2001).
Random forests.
\textit{Machine Learning}, 45(1), 5--32.

\bibitem{kruskal1952}
Kruskal, W.~H. \& Wallis, W.~A. (1952).
Use of ranks in one-criterion variance analysis.
\textit{Journal of the American Statistical Association}, 47(260), 583--621.

\bibitem{shapiro1965}
Shapiro, S.~S. \& Wilk, M.~B. (1965).
An analysis of variance test for normality (complete samples).
\textit{Biometrika}, 52(3/4), 591--611.

\bibitem{gray1984}
Gray, W.~M. (1984).
Atlantic seasonal hurricane frequency. Part I: El Ni\~no and 30~mb
quasi-biennial oscillation influences.
\textit{Monthly Weather Review}, 112(9), 1649--1668.

\bibitem{klotzbach2010}
Klotzbach, P.~J. (2010).
On the Madden--Julian oscillation--Atlantic hurricane relationship.
\textit{Journal of Climate}, 23(2), 282--293.

\bibitem{misra2009}
Misra, V. (2009).
\textit{Vedic Astrology: An Integrated Approach.}
Motilal Banarsidass Publishers.

\bibitem{emanuel2003}
Emanuel, K. (2003).
Tropical cyclones.
\textit{Annual Review of Earth and Planetary Sciences}, 31(1), 75--104.

\end{thebibliography}

% ══════════════════════════════════════════════════════════════════════════════
% Appendices
% ══════════════════════════════════════════════════════════════════════════════
\appendix

\section{Feature Block Definitions}
\label{app:blocks}

\begin{description}
  \item[Meteorological block (14 features)] \textit{wind}, \textit{pressure},
    \textit{latitude}, \textit{longitude}, \textit{wind\_change},
    \textit{pressure\_change}, \textit{displacement}, \textit{movement\_speed},
    \textit{hours\_elapsed}, \textit{spline\_bearing\_deg},
    \textit{spline\_curvature}, \textit{spline\_displacement},
    \textit{spline\_length\_deg}, \textit{elevation}.
  \item[Panchang block (18 features)] Tithi, Nakshatra, Yoga, Karana,
    Paksha, Moonsign, Sunsign, Drik Ritu, Drik Ayana, Samvatsara (all numeric),
    \textit{dinamana\_minutes}, \textit{ratrimana\_minutes}, day-of-week,
    day-of-year, month, hour\_utc, \textit{lahiri\_ayanamsha}.
  \item[Scientific block (190 features)] All Skyfield per-body position
    features (RA/Dec, ecliptic lon/lat, altitude/azimuth, sin/cos transforms)
    plus Vedic planetary numeric quantities (full degree, padam, speed, RA,
    Dec). \emph{Excludes} spline-prefixed features to avoid leakage.
\end{description}

\section{Pipeline Audit Summary}
\label{app:audit}

\begin{center}
\small
\begin{tabular}{llll}
  \toprule
  \textbf{ID} & \textbf{Script} & \textbf{Issue} & \textbf{Severity} \\
  \midrule
  A2 & 06 & Rahu $=$ Ketu retrograde (astronomical constraint) & Critical \\
  B1 & 02 & No multiple-comparison correction (fixed: BH FDR) & Medium \\
  C1 & 03 & Negative CV $R^2$; feature ranks unreliable & Critical \\
  C3 & 03 & Same-storm CV leakage (fixed: GroupKFold) & Medium \\
  D1 & 04 & $R^2$ discrepancy vs Stage 3 (documented) & Critical \\
  E2 & 05 & Circular degrees treated as linear (fixed: sin/cos) & Critical \\
  E3 & 05 & Retrograde: no $p$-values (fixed: Mann--Whitney $U$) & Medium \\
  F2 & 06 & Vedic model worse than baseline (flagged) & High \\
  F5 & 06 & Samvatsara ANOVA = interannual confound (documented) & High \\
  F6 & 06 & No Bonferroni correction (fixed: $\alpha_{\text{adj}} = 0.007$) & Medium \\
  G3 & 07 & Spline feature leakage (removed from sci.\ block) & Critical \\
  G2 & 07 & Best model = meteo\_only, not full (flagged) & High \\
  \bottomrule
\end{tabular}
\end{center}

\section{Consolidated Evidence Table}
\label{app:evidence}

\begin{longtable}{p{2cm}p{5.5cm}lp{4cm}}
  \caption{Full 16-finding evidence table with source scripts, effect sizes,
  and verdicts.}
  \label{tab:evidence_full} \\
  \toprule
  \textbf{ID} & \textbf{Finding} & \textbf{Source} & \textbf{Verdict} \\
  \midrule
  \endfirsthead
  \toprule
  \textbf{ID} & \textbf{Finding} & \textbf{Source} & \textbf{Verdict} \\
  \midrule
  \endhead
  \midrule \multicolumn{4}{r}{\textit{Continued on next page}} \\
  \endfoot
  \bottomrule
  \endlastfoot
  F01 & Vedic features degrade RMSE by 4.9\% (meteo+panchang) & 06, 07 & No predictive lift \\
  F02 & Panchang RMSE lift = 0.010 kt vs meteo 0.582 kt & 07 & $<$2\% explainable variance \\
  F03 & Lunar phase vs RI: $\chi^2=19.66$, $V=0.020$ & 06 & Stat.\ sig.; negligible $V$ \\
  F04 & Waning Crescent peak RI (1.33\%), Waxing Crescent low (0.79\%) & 06 & Real; lifecycle-dependent \\
  F05 & Yam speed vs bearing: $\rho=-0.178$, $r^2=3.2\%$ & 06 & Small; leakage risk \\
  F06 & Shani/Varun speed vs bearing: $\rho\approx-0.17$, $r^2\approx2.8\%$ & 06 & Same caveat as F05 \\
  F07 & Varun retrograde: wind $+4$ kt, $r_\text{rb}=-0.088$ & 05 & Small; retro 60\% of obs \\
  F08 & Shukra retrograde: wind $+5$ kt, $r_\text{rb}=-0.044$ & 05 & Small; genuinely rare (8\%) \\
  F09 & Rahu degree shift around RI onset: $F=18.75$, $\eta^2=0.097$ & 06 & Medium on $N=176$; replicate \\
  F10 & Mangal retrograde vs curvature: $F=19.34$, $\eta^2=0.0005$ & 06 & Trivially small \\
  F11 & Yoga KW $H=143$; 100\% strata sig & 02, 08 & Residual signal; small $\varepsilon^2$ \\
  F12 & Nakshatra KW $H=116$; 100\% strata sig & 02, 08 & Residual signal; small $\varepsilon^2$ \\
  F13 & Samvatsara ANOVA $F=65$, $p<10^{-130}$ & 06 & Confounded with calendar year \\
  F14 & Rahu degree vs $\Delta w$: $\rho=0.013$, $r^2=0.017\%$ & 06 & Negligible \\
  F15 & 4 Vedic features in SHAP/LIME top-14 consensus & 04 & Autocorrelated split inflates \\
  F16 & Meteo-only wins 4/5 forward-CV folds & 07 & Consistent across time \\
\end{longtable}

\end{document}